\section{Scope, status, and the three costs being reduced}
\label{sec:overview}

The objective is a mathematically justified acceleration of the ProveIt
unknot-recognition algorithm, with quasi-polynomial asymptotic
complexity as the research target. The repository already contains a
substantial collection of exact structural certificates, compressed group
operations, and relative Khovanov scanners. A useful continuation must
identify which expensive operations it actually avoids and which input-size
parameters remain uncontrolled. This report develops three such improvements,
and includes the code and negative experimental results needed to evaluate
them independently.

\begin{definition}[Input and target]
The primary input is a validated spherical planar diagram of a tame knot
with $n$ crossings and one component. Crossing and edge labels have
polynomially bounded encoding length after normalization. General
quasi-polynomial time means $\exp((\log(n+2))^{O(1)})$; several sufficient
conditions in this report yield the more specific bound
$\exp(O((\log(n+2))^2))=n^{O(\log n)}$. Bounds for a supplied compressed
presentation explicitly include its description length and provenance.
\end{definition}

Three different costs occur. The first is the number of real variables and
the polynomial degree in an exact representation query. The second is
cobordism-valued elimination performed before a decision certificate becomes
visible. The third is repeatedly eliminating the same graph interior while
preparing boundary observations at successive stages. Reducing any one of
these can be useful. None alone bounds the total number of intermediate
relative-complex objects or guarantees a short geometric decomposition.

\begin{table}[htbp]
\centering\small
\begin{tabularx}{\textwidth}{@{}p{0.23\textwidth}YY@{}}
\toprule
Result & Proved guarantee & Remaining condition or limitation\\
\midrule
Explicit knot-group presentation &
$\poly(N+r)(2+N/r)^{O(r)}$ bit operations &
Construct and verify the presentation; $N$ counts explicitly written letters.\\
Word-circuit cut &
$\poly(S+r+s)(d+2)^{O(r+t)}$ bit operations &
Pay for retained gates, formal degree, dense compilation, and provenance.\\
Pre-elimination first jet &
$t=N-2\rank A$ and $\kappa=N-\rank\bigl(\begin{smallmatrix}A&0\\B&A\end{smallmatrix}\bigr)$ &
Whole one-matching summand; general multiplier needs a knot completion.\\
Multiplier-one rigidity &
One matching preserves total rank under every classical link continuation &
Preserves a decision quantity; original quantum grading is not retained.\\
All-suffix responses &
$O(n(W+1)^2)$ field operations for every suffix &
Prescribed order; query count and relative-complex growth are separate.\\
Checkpoint scan &
$\poly(n)2^{O(g+w)}$ for its actual maximum gap and frontier &
No universal polylogarithmic bound on $g+w$ is proved.\\
\bottomrule
\end{tabularx}
\caption{The theorem boundaries. Letters reused inside individual sections
are defined locally; in the presentation rows $r$ counts generators and $N$
counts relator letters, while in the first-jet row $N$ counts complex objects.}
\label{tab:results}
\end{table}

\subsection{What is established and what is still research}

The presentation theorems are exact decision theorems using established
existential real algorithms. The delivered compiler emits exact integer
polynomials, but its optional practical solver is NLSAT through Z3; the
theoretical running-time bound is not attributed to that solver. The
first-jet scanner is a complete rank-decision method when its underlying
exact computation is allowed to finish, with a bounded optional observer
and an unchanged mathematical fallback. The streaming-response module is
a subroutine with a proved arithmetic bound, and is not installed in the
default recognition pipeline.

The development is based on the pinned revision \basecommit\ of the
ProveIt repository \cite{ProveItBaseline}. Prior reports establish the
binary Jones contraction, scalar residue and radical analysis, first-jet
knot-completion multiplier, and fixed-stage boundary responses. The new
work is the quantitative presentation compilation, the placement of the
first-jet test before elimination together with its multiplier-one link
extension and flat-origin constraint, and reuse of boundary elimination
across all changing suffixes. Established ingredients receive attribution;
the fact that a combination is proved here does not establish historical
priority for every ingredient or parameterized consequence.

% Optional introduction/footnote material, status checked 8 October 2026.
% Bibliography keys are supplied in su2_sources.bib.

\paragraph{Current literature and the meaning of the target.}
The present work does not claim a new general quasipolynomial algorithm.
Lackenby's March 2021 author slides announce a running time
$2^{O((\log n)^3)}=n^{O((\log n)^2)}$ and outline a hierarchy approach
\cite{Lackenby2021Slides}. We use that source's displayed bound.
His July 2026 preprint proves a new hierarchy
algorithm for incompressibility, with unknot recognition as a consequence,
and explicitly leaves possible speed-ups to future work
\cite{Lackenby2026Hierarchies}.  It should therefore not be cited as a
completed proof of the announced quasipolynomial bound.  The present
algebraic sufficient condition targets $n^{O(\log n)}$ on a specified
residual class, under explicitly stated and verified compression hypotheses.
These are different statements.

A September 2026 preprint by Musick claims a polynomial-time procedure
for producing locally minimal bridge presentations and, consequently,
unknot recognition in $P$ \cite{Musick2026}.  Version 2 is dated
20 September 2026.  That claim has not been independently established by
this study and none of the results or code here depends on it.  Recording
its existence avoids the inaccurate claim that no polynomial-time preprint
has appeared; treating it as an unverified recent claim avoids replacing
proof checking with an abstract-level assertion.  An independent audit
would have to establish the completeness of the permitted presentation
moves and the polynomial bounds on their compressed execution, as well as
the correctness of the initial diagram conversion.

\paragraph{What the hierarchy bounds suggest.}
Lackenby's 2026 Proposition 9.1 bounds the number of main iterations by
$L(g+1)^L$, where $L$ controls hierarchy length and $g$ controls the
relevant pattern complexity.  The later sections develop efficient
hierarchy certificates and compressed cutting operations
\cite{Lackenby2026Hierarchies}.  Purely as an implication of that bound,
polynomial $g$ and $L=O(\log n)$ would yield $n^{O(\log n)}$
iterations, while $L=O((\log n)^2)$ yields
$n^{O((\log n)^2)}$.  Counting iterations alone still requires a bound
on the bit cost of each iteration and on the construction of the
surfaces used.  The 2021 slides describe polynomial genus, compressed
hierarchies, multisurfaces, and Heegaard/Cheeger-region arguments as the
components of that route \cite{Lackenby2021Slides}.  A rigorous
implementation project would need exact data formats, verified transition
rules, and complete size and bit-complexity bounds for those components.


\subsection{A stronger sufficient condition is useful even with obstructions}

For a word circuit, let $r$ count original generators, let $t$ count
retained gates, and let $d$ bound the formal degrees after cutting. The
explicit target
\[
             (r+t)\log(d+2)=O((\log n)^2)
\]
turns an informal request to ``compress the group'' into a testable theorem
statement. It must be accompanied by polynomial construction and verification
costs. Connected sums of trefoils, and a prime pretzel family, have group
rank growing linearly with crossing count. Thus this condition cannot hold
universally for presentations of the entire knot group. The same knots
have inexpensive coloring obstructions, so they do not refute a portfolio
that first removes such cases. A meaningful residual-class theorem must
state those filters and decompositions exactly.

Likewise, the first-jet scan supplies an explicit target $g+w=O(\log^2 n)$
on the executed order, including temporary crossing allocation and the
last segment. The streaming response theorem reduces the cost of observing
such an order without proving that the order exists. These formulations
make a future quasi-polynomial claim susceptible to proof checking rather
than extrapolation from favorable timings.
